\documentclass[12pt,oneside]{book}

% =====================================================
% PAQUETES
% =====================================================
\usepackage[utf8]{inputenc}
\usepackage[T1]{fontenc}
\usepackage[spanish,es-nodecimaldot]{babel}

\usepackage[
    top=2.5cm,
    bottom=2.5cm,
    left=3cm,
    right=2.5cm,
    headheight=15pt
]{geometry}

\usepackage{amsmath}
\usepackage{amssymb}
\usepackage{mathtools}

\usepackage{graphicx}
\usepackage{float}
\usepackage{caption}
\usepackage{subcaption}

\usepackage{booktabs}
\usepackage{array}
\usepackage{longtable}
\usepackage{tabularx}

\usepackage{enumitem}

\usepackage{xcolor}
\usepackage[most]{tcolorbox}

\usepackage{fancyhdr}

\usepackage{ragged2e}

\graphicspath{
    {./img/}
    {./graficos/}
    {./figuras/}
}

% =====================================================
% CONFIGURACIÓN
% =====================================================
\setcounter{tocdepth}{2}

\setlength{\parindent}{0pt}
\setlength{\parskip}{0.5em}

\setlist[itemize]{
    topsep=0.3em,
    itemsep=0.2em
}
\setlist[enumerate]{
    topsep=0.3em,
    itemsep=0.2em
}

\clubpenalty=10000
\widowpenalty=10000

% =====================================================
% METADATOS
% =====================================================
\newcommand{\universidad}{Universidad de Buenos Aires (UBA)}
\newcommand{\materia}{Derecho Civil}
\newcommand{\titulo}{Régimen de Capacidad de las Personas Menores de Edad}
\newcommand{\autor}{Isaac Edgar Camacho Ocampo}
\newcommand{\anio}{2024}

% =====================================================
% COMANDOS MATEMÁTICOS
% =====================================================
\newcommand{\abs}[1]{\left|#1\right|}
\newcommand{\norm}[1]{\left\lVert#1\right\rVert}
\newcommand{\vect}[1]{\mathbf{#1}}
\newcommand{\deriv}[2]{\frac{d#1}{d#2}}
\newcommand{\pderiv}[2]{\frac{\partial#1}{\partial#2}}

% =====================================================
% ENTORNOS / CAJAS ACADÉMICAS
% =====================================================
\newtcolorbox{definition}[1]{
    enhanced,
    breakable,
    title={Definición: #1},
    fonttitle=\bfseries,
    colback=blue!3,
    colframe=blue!50!black,
    boxrule=0.8pt,
    arc=2mm
}

\newtcolorbox{important}{
    enhanced,
    breakable,
    title={Importante},
    fonttitle=\bfseries,
    colback=yellow!8,
    colframe=orange!70!black,
    boxrule=0.8pt,
    arc=2mm
}

\newtcolorbox{example}[1]{
    enhanced,
    breakable,
    title={Ejemplo: #1},
    fonttitle=\bfseries,
    colback=green!4,
    colframe=green!40!black,
    boxrule=0.8pt,
    arc=2mm
}

\newtcolorbox{formula}[1]{
    enhanced,
    breakable,
    title={Fórmula / Regla: #1},
    fonttitle=\bfseries,
    colback=purple!4,
    colframe=purple!50!black,
    boxrule=0.8pt,
    arc=2mm
}

\newtcolorbox{exercise}[1]{
    enhanced,
    breakable,
    title={Ejercicio: #1},
    fonttitle=\bfseries,
    colback=gray!5,
    colframe=gray!60!black,
    boxrule=0.8pt,
    arc=2mm
}

\newtcolorbox{note}{
    enhanced,
    breakable,
    title={Nota},
    fonttitle=\bfseries,
    colback=white,
    colframe=gray!50,
    boxrule=0.6pt,
    arc=2mm
}

% =====================================================
% ENCABEZADOS Y PIES
% =====================================================
\pagestyle{fancy}
\fancyhf{}

\fancyhead[L]{\nouppercase{\leftmark}}
\fancyhead[R]{\materia}

\fancyfoot[C]{\thepage}

\renewcommand{\headrulewidth}{0.4pt}
\renewcommand{\footrulewidth}{0pt}

\fancypagestyle{plain}{
    \fancyhf{}
    \fancyfoot[C]{\thepage}
    \renewcommand{\headrulewidth}{0pt}
}

% =====================================================
% HIPERVÍNCULOS Y METADATOS PDF
% =====================================================
\usepackage[
    colorlinks=true,
    linkcolor=blue!50!black,
    citecolor=green!40!black,
    urlcolor=blue!60!black,
    pdftitle={Régimen de Capacidad de las Personas Menores de Edad},
    pdfauthor={Isaac Edgar Camacho Ocampo},
    pdfsubject={Apuntes universitarios},
    bookmarks=true
]{hyperref}

% =====================================================
% DOCUMENTO
% =====================================================
\begin{document}

% -----------------------------------------------------
% PORTADA
% -----------------------------------------------------
\begin{titlepage}

\thispagestyle{empty}

\begin{center}

\vspace*{1cm}

{\large \textsc{\universidad}}

\vspace{3cm}

{\Huge \textbf{\materia}}

\vspace{1cm}

{\LARGE \titulo}

\vspace{0.8cm}

\textit{Comprender los fundamentos de cada instituto es el primer paso para transformar el estudio en criterio profesional.}

\vspace{1.2cm}

\rule{0.70\textwidth}{0.5pt}

\vspace{0.5cm}

{\large Apuntes de estudio}

\vspace{0.5cm}

\rule{0.70\textwidth}{0.5pt}

\vfill

{\large \autor}

\vspace{0.3cm}

{\large \anio}

\end{center}

\vfill

\begin{flushleft}
\textit{Este es un modesto aporte a los estudiantes de \textquotedblleft Derecho Civil\textquotedblright, no pretende sustituir el material oficial de la materia.}
\end{flushleft}

\end{titlepage}

% -----------------------------------------------------
% ÍNDICE
% -----------------------------------------------------
\tableofcontents

% =====================================================
% INTRODUCCIÓN Y RELEVANCIA DE LA UNIDAD
% =====================================================
\chapter*{Introducción a la unidad}
\addcontentsline{toc}{chapter}{Introducción a la unidad}
\markboth{Introducción a la unidad}{Introducción a la unidad}

Esta unidad aborda el régimen de capacidad de las personas menores de edad en el Código Civil y Comercial de la Nación (CCC), con especial atención a los principios rectores que lo estructuran, a los actos que la ley autoriza realizar según la edad, y a las consecuencias jurídicas y modos de finalización de la incapacidad de ejercicio.

Se trata de un contenido de aplicación directa en las áreas de familia, niñez y adolescencia, salud, contratos laborales y responsabilidad civil: determina qué puede hacer un menor por sí mismo, cuándo intervienen sus representantes y cuáles son los límites de esa representación, cuestiones relevantes para redactar contratos válidos, asesorar a familias, litigar en juzgados de familia o defender los derechos de niños y adolescentes en sede judicial.

Asimismo, la comprensión del régimen de capacidad progresiva permite explicar por qué ciertos actos realizados durante la minoría de edad tienen consecuencias jurídicas particulares, por qué los representantes legales intervienen en determinados contratos, o por qué un adolescente de trece años puede decidir sobre determinados tratamientos de salud. La materia atraviesa, además, debates actuales sobre autonomía, identidad y derechos de los adolescentes.

\begin{example}{El recorrido de la capacidad progresiva}
Los temas de esta unidad pueden entenderse como un recorrido de menor a mayor autonomía. Imaginemos a una persona, \textbf{Juan}, que nació sin capacidad de ejercicio y fue creciendo: a los 13 años pudo decidir sobre un tratamiento médico sencillo; a los 16 pudo trabajar, firmar contratos y cuidar su propio cuerpo como un adulto; al contraer matrimonio antes de los 18 años quedó emancipado, aunque con ciertas restricciones para proteger su patrimonio. Así se articula la tensión entre \textbf{autonomía y protección} que atraviesa toda esta unidad.
\end{example}

% =====================================================
% CAPÍTULO 1
% =====================================================
\chapter{Régimen general de capacidad}

\section{Concepto de menor de edad}

El Código Civil y Comercial de la Nación (CCC) denomina \textbf{personas menores de edad} a todas aquellas que no han cumplido los dieciocho años en la República Argentina.

\begin{definition}{Artículo 25 CCC — Menor de edad y adolescente}
Menor de edad es la persona que no ha cumplido 18 años. En caso de que hubiera cumplido 13 años, se lo denomina adolescente.
\end{definition}

\section{La incapacidad de ejercicio como medida de protección}

El régimen aplicable a las personas menores de edad se caracteriza, como punto de partida, por una \textbf{incapacidad de ejercicio}. Esta incapacidad no debe entenderse como una disminución de la persona: el menor de edad es, ante todo, un sujeto con dignidad, titular de derechos y con aptitud para adquirirlos, cuestión propia del campo de las relaciones básicas de la persona dentro del derecho civil.

Toda persona nace en el contexto de una familia, y por ello el ordenamiento confía primero en los padres como representantes del menor. Los progenitores ofrecen el ámbito de educación, crianza y contención necesario para el desarrollo de la persona: quienes dan la vida son, en principio, quienes se ocupan de criar y acompañar.

\begin{important}
La incapacidad de ejercicio \textbf{no es una sanción} ni una disminución de la persona. Es una herramienta jurídica que \textbf{protege} al menor, supliendo su falta de madurez con la actuación de sus representantes legales.
\end{important}

\section{Fundamento de la incapacidad: la inmadurez progresiva}

\begin{example}{El niño de diez años que hereda un inmueble}
Un niño o una niña de 10 años que hereda un inmueble no sabe el valor de las cosas, no sabe cómo negociar un contrato, no sabe bien qué hacer: quizás termina haciendo un negocio que lo perjudica, o incluso personas inescrupulosas pueden abusar de su inexperiencia para quedarse con sus bienes. Por ello se le designa un representante que actúe en su nombre.
\end{example}

El derecho contempla, a la vez, el \textbf{derecho a ser oído}: la relación entre padres e hijos presupone una manera de ir educando y llevando hacia la madurez a una persona que es necesariamente inmadura.

\subsection{La mayoría de edad a través del tiempo en Argentina}

\begin{table}[H]
\centering
\caption{Evolución histórica de la mayoría de edad en Argentina}
\begin{tabularx}{\textwidth}{Xcl}
\toprule
\textbf{Período} & \textbf{Edad de mayoría} & \textbf{Norma} \\
\midrule
Código Civil histórico & 22 años & Código de Vélez Sársfield \\
Reforma intermedia & 21 años & Modificación legislativa \\
Ley 26.579 (2009) & 18 años & Ley 26.579 \\
Código Civil y Comercial (2015) & 18 años & Ratificado por el CCC \\
\bottomrule
\end{tabularx}
\end{table}

La fijación de la mayoría de edad en los 18 años mantiene coherencia con el sistema de la Convención sobre los Derechos del Niño, que define como niño a todo ser humano hasta esa edad. Existe, entonces, \textbf{coherencia entre el Código Civil y la Convención de los Derechos del Niño}.

Corresponde advertir, además, sobre el riesgo de reducir la mayoría de edad con un espíritu meramente demagógico: bajar la mayoría de edad también implica dejar de brindar alimentos al joven. Por esa razón la ley estableció que, aunque se sea mayor de edad a los 18 años, en algunos casos debe continuarse con la provisión de alimentos, en tanto la persona todavía se encuentra en formación y no cuenta con un trabajo independiente.

\begin{note}
El régimen alimentario post-mayoría de edad extiende la obligación alimentaria hasta los 21 años con carácter general, y hasta los 25 años si el joven se encuentra en proceso de formación educativa o universitaria.
\end{note}

\section{La categoría de \textquotedblleft adolescente\textquotedblright}

La segunda parte del artículo 25 establece que, si el menor cumplió 13 años, se lo denomina \textbf{adolescente}. Esta categoría ha sido criticada por no constituir un concepto plenamente técnico desde el punto de vista jurídico, aunque se encuentra consagrada en el Código. El adolescente es, en consecuencia, el menor de entre 13 y 18 años.

Respecto de si el principio general es la capacidad o la incapacidad, existe debate doctrinario. Un sector de la doctrina sostiene que, por aplicación de las convenciones internacionales, el principio sería la capacidad. Otra parte relevante de la doctrina argentina entiende, en cambio, que la regla es la \textbf{incapacidad}, que va admitiendo una capacidad progresiva.

\begin{definition}{Artículo 26 CCC — Primera parte}
La persona menor de edad ejerce sus derechos a través de sus representantes legales. No obstante, la que cuenta con edad y grado de madurez suficiente puede ejercer por sí los actos que le son permitidos por el ordenamiento jurídico.
\end{definition}

\section*{Repaso — Cuadro Cornell}
\addcontentsline{toc}{section}{Repaso — Cuadro Cornell}

\begin{table}[H]
\centering
\begin{tabularx}{\textwidth}{
    >{\raggedright\arraybackslash}p{0.30\textwidth}
    >{\raggedright\arraybackslash}X
}
\toprule
\textbf{Preguntas / palabras clave} & \textbf{Notas / respuestas} \\
\midrule

\textbf{¿Quién es considerado menor de edad por el CCC?} &
Toda persona que no ha cumplido 18 años en la República Argentina (Art. 25 CCC). \\

\textbf{¿Por qué la incapacidad protege al menor?} &
Porque suple la falta de madurez del menor con la actuación de sus representantes legales, evitando que celebre negocios perjudiciales o sea víctima de abusos; no constituye una sanción ni una disminución de su dignidad. \\

\textbf{¿Qué diferencia al niño del adolescente?} &
El adolescente es la persona menor de edad que cumplió 13 años (segunda parte del Art. 25 CCC); por debajo de esa edad se lo considera niño. \\

\textbf{¿Por qué la ley fija una edad y no la madurez real?} &
Porque establecer un límite etario ofrece certeza y evita que cada caso deba discutirse individualmente; la ley presume una inmadurez progresiva que se va atenuando con el crecimiento. \\

\midrule
\multicolumn{2}{p{\textwidth}}{
\textbf{Resumen:} El CCC considera menor de edad a quien no cumplió 18 años, y adolescente a quien tiene entre 13 y 18. La incapacidad de ejercicio que rige a los menores es una medida de protección, no una sanción, fundada en la inmadurez progresiva propia del desarrollo de la persona.
} \\
\bottomrule
\end{tabularx}
\end{table}

% =====================================================
% CAPÍTULO 2
% =====================================================
\chapter{Principios rectores de la capacidad progresiva (Arts. 26 y 639 CCC)}

El artículo 639, ubicado en el capítulo de la responsabilidad parental (Libro Segundo, Derecho de Familia), enuncia tres principios fundamentales que rigen el ejercicio de los derechos de los niños, niñas y adolescentes.

\section{Interés superior del niño}

Este principio tiene como fuente la \textbf{Convención sobre los Derechos del Niño} de 1989, que la Argentina ratificó por Ley 23.849 e incorporó con jerarquía constitucional en la reforma constitucional de 1994, mediante el artículo 75, inciso 22.

La palabra \textquotedblleft interés\textquotedblright{} remite a un valor objetivo: qué es lo mejor para ese niño, cuál es el mayor bien para él.

\begin{definition}{Convención sobre los Derechos del Niño — Jerarquía constitucional}
Ratificada por Ley 23.849. Incorporada con jerarquía constitucional por la Reforma de 1994, Art. 75 inc. 22 de la Constitución Nacional.
\end{definition}

\section{Autonomía progresiva}

La autonomía es progresiva: a medida que el niño crece en madurez y en edad, va disminuyendo la representación de los padres y va aumentando la capacidad de los niños, niñas y adolescentes para ejercer por sí sus propios derechos.

Esta autonomía progresiva se encuentra empalmada con la regla del artículo 26: las personas menores de edad pueden ejercer aquellos derechos que la ley les autoriza. La autonomía va acompañada, entonces, de autorizaciones legales que el ordenamiento otorga progresivamente.

\section{Derecho a ser oído}

Este principio proviene del artículo 12 de la Convención sobre los Derechos del Niño y tiene su correlato en el artículo 26 del CCC: el derecho del niño a ser oído en todos los asuntos de su vida y a participar en las decisiones fundamentales sobre la conducción de su persona.

\begin{example}{Divorcio con hijos menores}
En un divorcio con hijos menores, corresponde escuchar a los niños respecto de las cuestiones que los afectan directamente, como la decisión sobre dónde vivirán o a qué colegio asistirán.
\end{example}

\section{Representación de los menores de edad}

La representación de los menores de edad corresponde, en primer lugar, a sus \textbf{padres}. A falta de los padres —por fallecimiento, privación de la responsabilidad parental u otra causa— la representación recae en los \textbf{tutores} (Art. 101 CCC).

Adicionalmente, la Ley 26.061 incorporó la figura del \textbf{abogado del niño}: un letrado que lo asesora cuando debe estar en juicio, especialmente ante un conflicto de intereses que no puede ser representado por sus progenitores.

Finalmente, el \textbf{Ministerio Público} ofrece una representación genérica —judicial o extrajudicial, principal o complementaria— regulada en el artículo 103 del CCC.

\begin{table}[H]
\centering
\caption{Sistema de representación del menor de edad}
\begin{tabularx}{\textwidth}{Xll}
\toprule
\textbf{Representante} & \textbf{Fuente legal} & \textbf{Carácter} \\
\midrule
Padres / Progenitores & Art. 101 CCC — Resp. parental & Principal y primario \\
Tutores & Art. 101 CCC & Supletorio (a falta de padres) \\
Abogado del niño & Ley 26.061 & Técnico — en juicio \\
Ministerio Público & Art. 103 CCC & Genérico — judicial y extrajudicial \\
\bottomrule
\end{tabularx}
\end{table}

\section{Determinación del grado de madurez suficiente}

La ley, en ocasiones, establece un límite de edad fijo (por ejemplo, 10, 13 o 16 años); en otros casos utiliza una fórmula amplia: \textquotedblleft grado de madurez suficiente\textquotedblright. Esta fórmula ha recibido críticas de parte de la doctrina.

En principio, su determinación depende del ámbito en que se aplique. Cuando la cuestión se presenta ante un juez, será este —con el asesoramiento de su equipo interdisciplinario— quien determine si el menor cuenta con el grado de madurez suficiente exigido por el Código.

% TODO: contenido incompleto o ambiguo en las notas originales — no se identifica con precisión la obra o publicación en la que el autor citado (Rivero) desarrolla su crítica.
Cierto sector doctrinario, representado por el autor Rivero, critica esta indeterminación y sostiene que debería eliminarse en favor de edades fijas. Existe, según las notas originales, un proyecto de reforma del CCC —elaborado por una comisión presidida por dicho autor— que tiende a eliminar esa indeterminación.

\section*{Repaso — Cuadro Cornell}
\addcontentsline{toc}{section}{Repaso — Cuadro Cornell}

\begin{table}[H]
\centering
\begin{tabularx}{\textwidth}{
    >{\raggedright\arraybackslash}p{0.30\textwidth}
    >{\raggedright\arraybackslash}X
}
\toprule
\textbf{Preguntas / palabras clave} & \textbf{Notas / respuestas} \\
\midrule

\textbf{¿Qué significa el \textquotedblleft interés superior\textquotedblright{} en la práctica?} &
Es un valor objetivo: lo que resulta mejor para ese niño en particular, fuente de la Convención sobre los Derechos del Niño (1989), con jerarquía constitucional desde 1994 (Art. 75 inc. 22 CN). \\

\textbf{¿Cómo crece la autonomía con la edad?} &
A mayor madurez y edad, disminuye la representación de los padres y aumenta la capacidad del niño, niña o adolescente para ejercer por sí sus derechos, siempre dentro de las autorizaciones legales del Art. 26 CCC. \\

\textbf{¿En qué casos el juez debe escuchar al menor?} &
En todos los asuntos que afecten su vida y en las decisiones fundamentales sobre la conducción de su persona (Art. 12 CDN y Art. 26 CCC); por ejemplo, en un divorcio, respecto de dónde vivirá o a qué colegio asistirá. \\

\textbf{¿Quién actúa en nombre del menor?} &
En primer lugar los padres; a falta de ellos, los tutores (Art. 101 CCC); también el abogado del niño (Ley 26.061) en juicio, y el Ministerio Público con representación genérica (Art. 103 CCC). \\

\midrule
\multicolumn{2}{p{\textwidth}}{
\textbf{Resumen:} Los principios rectores del sistema —interés superior del niño, autonomía progresiva y derecho a ser oído— provienen de la Convención sobre los Derechos del Niño e integran el bloque constitucional. La representación del menor recae principalmente en los padres, y de modo supletorio o complementario en tutores, abogado del niño y Ministerio Público.
} \\
\bottomrule
\end{tabularx}
\end{table}

% =====================================================
% CAPÍTULO 3
% =====================================================
\chapter{Actos que la ley autoriza realizar según la edad}

\section{Tratamientos médicos (Art. 26)}

\begin{definition}{Artículo 26 CCC — Tratamientos médicos}
Desde los 13 años se presume que la persona menor de edad posee aptitud para decidir sobre tratamientos no invasivos o que no comprometan su vida o su salud. Si el tratamiento es invasivo y compromete su salud, se requiere la asistencia (consentimiento) tanto de la persona menor de edad como de sus padres. En caso de discordancia, se decide con base en la opinión médica, teniendo en cuenta el interés superior. Desde los 16 años, el adolescente es considerado como un adulto para el cuidado de su propio cuerpo.
\end{definition}

\begin{table}[H]
\centering
\caption{Régimen de decisiones médicas según la edad}
\begin{tabularx}{\textwidth}{lXX}
\toprule
\textbf{Edad} & \textbf{Tipo de tratamiento} & \textbf{¿Quién decide?} \\
\midrule
0 a 12 años & Cualquier tratamiento & Los padres (o tutores) \\
13 a 15 años & No invasivo / sin riesgo de vida & El propio menor (presunción de capacidad) \\
13 a 15 años & Invasivo / con riesgo de vida & Menor + padres (asistencia conjunta) \\
16 años en adelante & Cuidado del propio cuerpo & El adolescente como adulto \\
\bottomrule
\end{tabularx}
\end{table}

Respecto de los conflictos entre el adolescente y sus progenitores en materia de tratamientos invasivos, el Código establece que la discordancia se resuelve teniendo en cuenta el interés superior del menor y la opinión médica.

\section{Actos permitidos — Resumen por edad}

\subsection{Desde el nacimiento o desde temprana edad}

\begin{example}{Pequeños contratos — Art. 684}
Cuando un menor va al kiosco o a la librería del barrio y compra algo, está celebrando un contrato que no es nulo. El artículo 684 presume que el menor goza de capacidad para los pequeños contratos de menor cuantía.
\end{example}

\subsection{Desde los 13 años}

\begin{itemize}
    \item Decidir sobre tratamientos médicos no invasivos (Art. 26).
    \item Iniciar una acción para conocer sus orígenes (Art. 596).
    \item Reconocer hijos.
    \item Ejercer la responsabilidad parental, con la supervisión especial del artículo 644: los abuelos (progenitores de los adolescentes-padres) pueden intervenir en las cuestiones más trascendentes ante la inacción de los progenitores adolescentes, y pueden vetar decisiones contrarias al interés del niño.
    \item Participar en forma autónoma en juicio con autorización de sus progenitores (Art. 677).
\end{itemize}

\begin{definition}{Artículo 644 CCC — Responsabilidad parental de progenitores adolescentes}
Los progenitores adolescentes que tengan hijos menores ejercen la responsabilidad parental conforme a las reglas especiales del artículo 644, que establece una supervisión de sus propios progenitores (los abuelos del niño) en las decisiones más trascendentes.
\end{definition}

\subsection{Desde los 16 años}

\begin{itemize}
    \item Ser equiparado a los adultos en relación con el cuidado del propio cuerpo (Art. 26, in fine).
    \item Celebrar contrato de trabajo: la Ley 26.390 regula el trabajo infantil y prohíbe el trabajo de menores de 16 años en sus distintas formas. Desde los 16, se presume que pueden celebrar el contrato de trabajo solos si viven de forma independiente.
    \item Ser donante de sangre con autorización de sus padres.
\end{itemize}

\begin{definition}{Ley 26.390 — Trabajo infantil}
Prohíbe el trabajo de personas menores de 16 años en sus distintas formas. Desde los 16 años se presume que el adolescente puede celebrar un contrato de trabajo con la autorización de sus padres, o bien solo si vive de forma independiente. Esta norma está en clave con el Convenio de la Organización Internacional del Trabajo.
\end{definition}

\section{Caso especial: el título habilitante (Art. 30)}

El artículo 30 regula el supuesto de la persona menor de edad que adquiere un título habilitante. Por ejemplo, quien obtuvo el título de profesor de piano en un conservatorio a los 16 años puede ejercer esa profesión y administrar los fondos que gane.

La doctrina discute qué ocurre si el menor tiene el título habilitante pero es menor de 16 años: un sector entiende que el artículo 30 no distingue edad, de modo que, si tiene título habilitante, puede ejercer; otro sector señala que, si pudo obtener el título —que requiere validación del Ministerio de Educación y requisitos formales—, ya demostró la madurez suficiente para ejercer esa profesión.

\begin{definition}{Artículo 30 CCC — Título habilitante}
La persona menor de edad que obtiene un título habilitante de profesión puede ejercerla por cuenta propia, celebrar todos los actos que la profesión le exija y administrar y disponer libremente de los bienes que adquiera con el producto de su trabajo.
\end{definition}

\section{El matrimonio de personas menores de edad}

\begin{table}[H]
\centering
\caption{Régimen matrimonial según la edad}
\begin{tabularx}{\textwidth}{lcX}
\toprule
\textbf{Edad} & \textbf{¿Puede contraer matrimonio?} & \textbf{Requisito} \\
\midrule
18 años o más & Sí & Ninguno adicional — regla general \\
Entre 16 y 17 años & Sí & Autorización de representantes legales; si la niegan, dispensa judicial \\
Menor de 16 años & Eventualmente sí & Siempre con dispensa judicial \\
\bottomrule
\end{tabularx}
\end{table}

\section*{Repaso — Cuadro Cornell}
\addcontentsline{toc}{section}{Repaso — Cuadro Cornell}

\begin{table}[H]
\centering
\begin{tabularx}{\textwidth}{
    >{\raggedright\arraybackslash}p{0.30\textwidth}
    >{\raggedright\arraybackslash}X
}
\toprule
\textbf{Preguntas / palabras clave} & \textbf{Notas / respuestas} \\
\midrule

\textbf{¿Puede un menor decidir solo sobre su salud?} &
Desde los 13 años, sobre tratamientos no invasivos. Para tratamientos invasivos se requiere su consentimiento y el de sus padres. Desde los 16 años se lo equipara a un adulto en el cuidado de su propio cuerpo (Art. 26 CCC). \\

\textbf{¿Es válida la compra en el kiosco hecha por un menor?} &
Sí; el Art. 684 presume que el menor tiene capacidad para celebrar pequeños contratos de menor cuantía. \\

\textbf{¿Puede trabajar un adolescente de 16 años?} &
Sí; la Ley 26.390 prohíbe el trabajo de menores de 16 años, y desde esa edad se presume que el adolescente puede celebrar contrato de trabajo, en especial si vive de forma independiente. \\

\textbf{¿A qué edad se puede contraer matrimonio?} &
Desde los 18 años sin requisitos adicionales; entre 16 y 17 con autorización de los representantes o dispensa judicial; por debajo de los 16, siempre con dispensa judicial. \\

\textbf{¿Qué pasa si un menor obtiene un título profesional?} &
Puede ejercer la profesión por cuenta propia y administrar los bienes obtenidos con su trabajo (Art. 30 CCC); la doctrina discute su aplicación cuando el título se obtiene antes de los 16 años. \\

\midrule
\multicolumn{2}{p{\textwidth}}{
\textbf{Resumen:} La ley reconoce capacidades específicas según la edad: desde los 13 años para decisiones médicas simples y ciertos actos de estado civil; desde los 16 para el trabajo y el cuidado autónomo del propio cuerpo; y desde el nacimiento para los pequeños contratos de menor cuantía. El matrimonio antes de los 18 años requiere autorización o dispensa judicial según la edad.
} \\
\bottomrule
\end{tabularx}
\end{table}

% =====================================================
% CAPÍTULO 4
% =====================================================
\chapter{Consecuencias jurídicas y fin de la incapacidad}

\section{Nulidad relativa de los actos celebrados sin representación}

\begin{example}{El cuadro heredado}
Una persona de 13 años heredó un cuadro valorado en 2.000 dólares. Aburrida durante la cuarentena, lo puso a la venta en una plataforma \textit{online} y obtuvo 4.000 dólares, el doble de su valor. Sus padres se enteraron y, en lugar de anular el acto, decidieron confirmarlo porque había resultado un buen negocio. El acto era nulo de nulidad relativa, pero los padres pudieron confirmarlo en beneficio del menor.
\end{example}

La consecuencia de que no se respete el régimen de capacidad y el menor actúe sin representación es la nulidad de carácter relativo. La nulidad es relativa porque:

\begin{itemize}
    \item Está puesta para proteger al menor, no para perjudicarlo.
    \item Si el acto lo beneficia, los padres pueden confirmarlo.
    \item Solo el menor (o sus representantes) puede alegar esa nulidad, no la contraparte.
\end{itemize}

\section{Responsabilidad civil por daños causados por menores}

\begin{table}[H]
\centering
\caption{Responsabilidad por daños según la edad del menor}
\begin{tabularx}{\textwidth}{Xcc}
\toprule
\textbf{Situación} & \textbf{Responde el menor} & \textbf{Responden los padres} \\
\midrule
Menor de 10 años — acto ilícito & No (acto reputado involuntario) & Sí, siempre (responsabilidad objetiva) \\
Mayor de 10 años — acto ilícito & Sí, con su propio patrimonio & Sí, también (responsabilidad objetiva) \\
\bottomrule
\end{tabularx}
\end{table}

Los padres responden siempre por los daños causados por sus hijos menores. Se trata de una responsabilidad objetiva —fundada en el deber de vigilancia— que no cesa aunque el menor también responda con su propio patrimonio cuando tiene más de 10 años.

\section{Fin de la incapacidad: la mayoría de edad}

\begin{important}
El día en que la persona cumple 18 años, se la considera automáticamente mayor de edad y plenamente capaz. No se requiere ningún acto formal: no hay que acudir al escribano ni al registro civil; basta con acreditar que se han cumplido los 18 años. Esto rige en el campo del derecho civil (contratos, familia, patrimonio) y no debe confundirse con la responsabilidad penal ni con el sufragio.
\end{important}

\section{Fin de la incapacidad: la emancipación por matrimonio}

La emancipación es la institución por la cual se confiere de forma anticipada la plena capacidad a una persona menor de edad, por alguno de los supuestos previstos por la ley. El CCC mantiene un único tipo de emancipación.

Históricamente existían dos formas: la emancipación por habilitación de edad —que operaba entre los 18 y los 21 años, cuando la mayoría se alcanzaba a los 21— y la emancipación por matrimonio. Cuando la mayoría de edad bajó a los 18 años, desapareció el instituto de la emancipación dativa o por habilitación de edad, que el nuevo Código no recoge. Queda, entonces, solo la emancipación por matrimonio.

\begin{definition}{Artículo 27 CCC — Emancipación}
La celebración del matrimonio antes de los 18 años emancipa a la persona menor de edad. La persona emancipada goza de plena capacidad de ejercicio con las limitaciones previstas en este Código. La emancipación es irrevocable. La nulidad del matrimonio no deja sin efecto la emancipación, excepto respecto del cónyuge de mala fe para quien el matrimonio fue anulado.
\end{definition}

\section{Restricciones a la capacidad del emancipado (Arts. 28 y 29)}

La emancipación no confiere plena capacidad sin restricciones. Los artículos 28 y 29 establecen limitaciones específicas.

\subsection{Artículo 28 — Actos prohibidos incluso con autorización judicial}

\begin{definition}{Artículo 28 CCC — Actos prohibidos a las personas emancipadas}
Las personas emancipadas no pueden, ni con autorización judicial:
\begin{enumerate}
    \item Aprobar las cuentas de sus tutores ni darles finiquito.
    \item Hacer donación de bienes que hubieran recibido a título gratuito.
    \item Ser fiadores.
\end{enumerate}
\end{definition}

El fundamento de cada prohibición puede resumirse así:

\begin{enumerate}
    \item \textbf{Aprobar cuentas de tutores}: el menor que tuvo tutor no puede aprobar las cuentas de su administración sino hasta ser mayor de edad; el Código sigue la regla del artículo 130.
    \item \textbf{Donación de bienes recibidos a título gratuito}: el emancipado no puede regalar lo que heredó. Esta norma proviene del Código de Vélez (1871) y no se modificó en 2015; busca evitar que el emancipado disponga de bienes cuyo valor todavía no comprende cabalmente.
    \item \textbf{Ser fiador}: la prohibición tiende a custodiar el patrimonio del emancipado.
\end{enumerate}

\subsection{Artículo 29 — Actos que requieren autorización judicial}

\begin{definition}{Artículo 29 CCC — Disposición de bienes a título oneroso}
El emancipado puede disponer de los bienes recibidos a título gratuito (venderlos, no donarlos), pero requiere autorización judicial. Dicha autorización solo procede en casos de absoluta necesidad o ventaja evidente.
\end{definition}

\begin{example}{El inmueble heredado por el emancipado}
El menor emancipado que heredó un inmueble no puede venderlo libremente: solo puede hacerlo con autorización judicial, y únicamente si existe una absoluta necesidad (como una operación médica urgente) o una ventaja evidente (un gran negocio). El CCC unificó aquí el sistema, eliminando la excepción del Código de Vélez que admitía prescindir de la dispensa cuando uno de los cónyuges era mayor de edad.
\end{example}

\begin{table}[H]
\centering
\caption{Cuadro comparativo: restricciones del emancipado}
\begin{tabularx}{\textwidth}{Xlcl}
\toprule
\textbf{Acto} & \textbf{Artículo} & \textbf{¿Puede hacerlo?} & \textbf{Condición} \\
\midrule
Aprobar cuentas del tutor & Art. 28 & No & Nunca, ni con autorización judicial \\
Donar bienes recibidos gratuitamente & Art. 28 & No & Nunca, ni con autorización judicial \\
Ser fiador & Art. 28 & No & Nunca, ni con autorización judicial \\
Vender bienes recibidos gratuitamente & Art. 29 & Sí, con límites & Requiere autorización judicial: necesidad absoluta o ventaja evidente \\
\bottomrule
\end{tabularx}
\end{table}

\section{Ley 27.364 — Mayoría de edad anticipada para egresados de institutos}

La Ley 27.364 estableció una mayoría de edad anticipada para los adolescentes que egresan de institutos y que no cuentan con cuidados parentales.

\begin{note}
Esta ley recibió críticas de la doctrina: se señaló que, precisamente por tratarse de adolescentes sin cuidados parentales y provenientes de condiciones de institucionalización, la mayoría de edad anticipada podría funcionar como una forma de abandono, aunque la ley contempla un programa estatal de protección.
\end{note}

\section*{Repaso — Cuadro Cornell}
\addcontentsline{toc}{section}{Repaso — Cuadro Cornell}

\begin{table}[H]
\centering
\begin{tabularx}{\textwidth}{
    >{\raggedright\arraybackslash}p{0.30\textwidth}
    >{\raggedright\arraybackslash}X
}
\toprule
\textbf{Preguntas / palabras clave} & \textbf{Notas / respuestas} \\
\midrule

\textbf{¿Qué pasa si un menor celebra un contrato sin representación?} &
El acto es nulo, pero de nulidad relativa: está puesta para proteger al menor, por lo que sus representantes pueden confirmarlo si le resulta beneficioso, y solo el menor o sus representantes pueden alegar la nulidad. \\

\textbf{¿Responden los padres por daños causados por sus hijos?} &
Sí, siempre, por responsabilidad objetiva fundada en el deber de vigilancia. Si el menor tiene más de 10 años, responde además con su propio patrimonio. \\

\textbf{¿Cómo y cuándo termina la incapacidad?} &
Principalmente por mayoría de edad, en forma automática el día del 18° cumpleaños, sin acto formal. Alternativamente, por emancipación mediante matrimonio (Art. 27 CCC). \\

\textbf{¿Qué efectos jurídicos produce el matrimonio de un menor?} &
Lo emancipa, otorgándole plena capacidad de ejercicio con las limitaciones de los Arts. 28 y 29 CCC. La emancipación es irrevocable, incluso si el matrimonio es declarado nulo (salvo para el cónyuge de mala fe). \\

\textbf{¿Qué actos no puede hacer quien se emancipó por matrimonio?} &
No puede aprobar cuentas de su tutor, donar bienes recibidos gratuitamente ni ser fiador (Art. 28), ni con autorización judicial. Puede vender bienes recibidos gratuitamente solo con autorización judicial, por necesidad absoluta o ventaja evidente (Art. 29). \\

\midrule
\multicolumn{2}{p{\textwidth}}{
\textbf{Resumen:} Los actos celebrados por un menor sin representación son nulos de nulidad relativa y pueden confirmarse si lo benefician. La incapacidad cesa automáticamente a los 18 años o, anticipadamente, por emancipación mediante matrimonio, institución que otorga plena capacidad pero con restricciones específicas en los artículos 28 y 29 del CCC.
} \\
\bottomrule
\end{tabularx}
\end{table}

% =====================================================
% CAPÍTULO FINAL — CORNELL INTEGRADOR
% =====================================================
\chapter{Repaso general — Cuadro Cornell integrador}

El siguiente cuadro sintetiza, en formato Cornell, los conceptos centrales de toda la unidad sobre régimen de capacidad de las personas menores de edad.

\begin{longtable}{
    >{\raggedright\arraybackslash}p{0.30\textwidth}
    >{\raggedright\arraybackslash}p{0.62\textwidth}
}
\caption{Cuadro Cornell integrador — Régimen de capacidad de las personas menores de edad} \\
\toprule
\textbf{Pregunta / concepto clave} & \textbf{Notas / respuesta / explicación} \\
\midrule
\endfirsthead

\multicolumn{2}{c}{\textit{(continuación)}} \\
\toprule
\textbf{Pregunta / concepto clave} & \textbf{Notas / respuesta / explicación} \\
\midrule
\endhead

\midrule
\multicolumn{2}{r}{\textit{continúa en la página siguiente}} \\
\endfoot

\bottomrule
\endlastfoot

\textbf{¿Qué es un menor de edad?} &
Toda persona que no ha cumplido 18 años (Art. 25 CCC). Si tiene entre 13 y 17 años, se lo denomina específicamente adolescente. \\

\textbf{¿Cuál es la regla sobre su capacidad?} &
La regla es la incapacidad de ejercicio. Esta incapacidad no es una sanción, sino una medida de protección que se suple con la representación de sus padres o tutores. \\

\textbf{¿Cuáles son los tres principios rectores?} &
1) Interés superior del niño (CDN, Ley 23.849, Art. 75 inc. 22 CN). 2) Autonomía progresiva: a mayor madurez y edad, mayor capacidad propia y menor representación parental. 3) Derecho a ser oído en todos los asuntos de su vida. \\

\textbf{¿Qué puede decidir un menor de 13 años sobre su salud?} &
Puede decidir solo sobre tratamientos médicos no invasivos o que no comprometan su salud o su vida. Para tratamientos invasivos se requiere su consentimiento y el de sus padres. Desde los 16 años se lo considera adulto en todo lo relativo al cuidado de su propio cuerpo. \\

\textbf{¿Qué son los pequeños contratos?} &
Los contratos de menor cuantía que cualquier menor puede celebrar, como comprar en un kiosco o una librería. El Art. 684 presume que el menor tiene capacidad para estos actos. Son plenamente válidos. \\

\textbf{¿Desde qué edad puede trabajar?} &
Desde los 16 años, conforme a la Ley 26.390, que prohíbe el trabajo de menores de esa edad. Se presume que el adolescente de 16 puede celebrar contrato de trabajo, especialmente si vive de forma independiente. \\

\textbf{¿Qué es la nulidad relativa?} &
El acto celebrado por un menor sin representación es nulo, pero de nulidad relativa: está puesta para protegerlo. Si el acto le beneficia, sus representantes legales pueden confirmarlo y darle plena validez. \\

\textbf{¿Cuándo termina la incapacidad?} &
Principalmente por mayoría de edad, el día del 18° cumpleaños, en forma automática, sin acto formal. Alternativamente, por emancipación mediante matrimonio (Art. 27 CCC), que confiere capacidad anticipada con ciertas restricciones. \\

\textbf{¿Qué restricciones tiene el emancipado?} &
Art. 28: no puede aprobar cuentas de su tutor, ni donar bienes recibidos gratuitamente, ni ser fiador, aun con autorización judicial. Art. 29: puede disponer a título oneroso de bienes gratuitos (venderlos), pero requiere autorización judicial por necesidad absoluta o ventaja evidente. \\

\textbf{Tensión fundamental de la materia} &
Toda la regulación de la capacidad de los menores está atravesada por la tensión entre autonomía (respeto por su voluntad y dignidad) y protección (resguardo de quien aún no ha alcanzado la madurez suficiente). La ley intenta equilibrarlas a través de la capacidad progresiva. \\

\midrule
\multicolumn{2}{p{0.94\textwidth}}{
\textbf{Síntesis final:} El régimen de capacidad de las personas menores de edad en el Código Civil y Comercial argentino parte de un principio claro: toda persona menor de 18 años es incapaz de ejercicio, no como sanción, sino como medida de protección. Esta incapacidad se complementa con un sistema de representación (padres, tutores, abogado del niño, Ministerio Público) y se va atenuando progresivamente a medida que el menor crece y madura. Los principios rectores del sistema —interés superior del niño, autonomía progresiva y derecho a ser oído— provienen de la Convención sobre los Derechos del Niño (1989) e integran el bloque constitucional desde 1994. La ley reconoce capacidades específicas según la edad: desde los 13 años para decisiones médicas simples y ciertos actos de estado civil; desde los 16 años para el trabajo y para el cuidado autónomo del propio cuerpo. Los actos celebrados sin representación son nulos de nulidad relativa y pueden ser confirmados si benefician al menor. La incapacidad cesa automáticamente a los 18 años o, anticipadamente, por emancipación mediante matrimonio —institución que confiere plena capacidad pero con restricciones específicas en los artículos 28 y 29 del CCC. La tensión permanente entre autonomía y protección es el hilo conductor de toda la materia.
} \\

\end{longtable}

\end{document}
